\section[E]{Subdirectories and GAMS Programs}\label{programs}

\begin{longtable}{l|l|p{4in}}
\texttt{\textbf{build/}}&&Working directory for dataset construction.\\
\hline
&       \texttt{build.gms} &Command script to read, filter and aggregate a set of GTAP datasets.
                        Default configuration: produces datasets \texttt{gtapingams.gdx}, \texttt{g20.gdx},
                        \texttt{g20\_macro.gdx} and \texttt{g20\_iea.gdx} for 2011.  Calls \texttt{code/flex2gdx.gms},
                        \texttt{code/filter.gms}, \texttt{code/gtapaggr.gms} and \texttt{code/cdecalib.gms}.\\[0.1in]
&       \texttt{filterchk.gms}& Command script to generate data for Table 6: Filtering results.
Calls \texttt{code/flex2gdx.gms}, \texttt{code/filter.gms}.  Results
are written to a PivotData sheet in \texttt{filterchk.xlsx}.\\
&       \texttt{test.gms}&               Command script which calls \texttt{bmkchk.gms}
                        for dataset \texttt{g20\_macro}, \texttt{mcpmge.gms} for
                        dataset \texttt{g20} and \texttt{cdechk.gms} for dataset
                        \texttt{g20}.\\[0.1in]
&       \texttt{bmkchk.gms}&    Benchmark consistency is a necessary, but not a
                        sufficient condition that the model is properly specified.
                        This routine checks for the \texttt{g20\_macro} dataset in both \MGE
                        and \MCP formats.  This routine produces a replication check
                        for the \GMR model and all the single region \SOE  models. \\[0.1in]
&       \texttt{mcpmge.gms}&    Is designed to perform a check of consistency for the
                        \MCP and \MGE models at a point away from the benchmark
                        equilibrium.  It can be used to verify that
                        a solution computed with \texttt{mge.gms} also solves \texttt{mcp.gms},
                        and that a solution computed with \texttt{mcp.gms} solves \texttt{mge.gms}.
                        The \GAMS savepoint and loadpoint commands are used for
                        this purpose.\\[0.1in]
&       \texttt{cdechk.gms}&    Verifies that the CDE demand system reproduces the
                        exogenous own-price and income demand elasticities
                        at the benchmark point.\\
\hline \multicolumn{3}{c}{}
\end{longtable}

\begin{longtable}{l|l|p{4in}}
\texttt{\textbf{code/}}&&               Code repository -- not a working directory.\\
\hline
&       \texttt{filter.gms}&    Filter routine, based on environment parameter --nd, the
                        number of decimal points in the filter.\\[0.1in]
&       \texttt{gtap12data.gms}& Utility routine for read a \GTAPINGAMS version 12 dataset.\\[0.1in]
&       \texttt{gtapaggr.gms}&  Dataset aggregation routine (call to \texttt{cdecalib.gms} following
                        aggregation must be made by user).\\[0.1in]
&       \texttt{aggr.gms}&      Utility routines called by gtapaggr.gms.\\
&       \texttt{chktarget.gms}&\\
&       \texttt{checkset.gms}&\\[0.1in]
&       \texttt{domain.gms}&    Standard purpose libinclude routine for extracting the nonzero domain of a parameter, using the mysterious
                        "\texttt{option pd\_1<p\_1;}" syntax.\\[0.1in]
&       \texttt{cdecalib.gms}&  Utility routine to recalibrate \CDE demand following dataset aggregation.\\[0.1in]
&       \texttt{mcp.gms}&               The canonical static \GTAPINGAMS model in algebraic format.\\[0.1in]
&       \texttt{mge.gms}&               The canonical static \GTAPINGAMS model in tabular \MPSGE format.\\[0.1in]
&       \texttt{gdpcalc.gms}&   Utility routine for reporting GDP on the basis of income,
                        final demand or sectoral value-added, compatible with
                        either the \texttt{mcp.gms} or \texttt{mge.gms} models.  First call
                        declares report parameters.  \\[0.1in]
&       \texttt{loadmdl.gms}&   Utility routines for switching demand systems (between
                        Cobb-Douglas, \LES and \CDE) and model closures (\GMR versus
                        \SOE).  The macros included here are described below.\\
\hline \multicolumn{3}{c}{}
\end{longtable}

\begin{longtable}{l|l|p{4in}}
\texttt{\textbf{forensics/}}&&  Produces some calculations assessing the economic consequences of
                proportional unilateral changes in trade taxes.\\
\hline \\
& \texttt{run.gms}& This \GAMS script which execute
\texttt{ssagen.gms}, \texttt{ssasolve.gms}  and \texttt{ssamerge.gms}.  \\
& \texttt{ssagen.gms}& This \GAMS script produces
\texttt{ssasolve.gms}  and \texttt{ssamerge.gms}\\
& \texttt{ssasolve.gms}& is a \GAMS programs
written by \texttt{ssagen.gms} which processes
\texttt{tgrid.gms} for several alternative model closures and demand
system specifications.\\
& \texttt{ssamerge.gms}& is a \GAMS programs written by
\texttt{ssagen.gms} which combines model results into a format
suitable for generating an Excel pivot report.\\
& \texttt{tgrid.gms}& Performs a sequence of counterfactual
calculations for a single region, with proportional scaling tariff and
export subsidy rates to average values between 0 and 10\%.  These
counterfactual calculations are computed with Cobb-Douglas, \LES and
\CDE demand.  Macro results (welfare, model status, and global
welfare) are returned in parameter \texttt{results(dsys,tlvl,*)}.  GDP
calculations are stored in
\texttt{gdpresults(dsys,tlvl,gdpcat,r,gdpitem)}. \\
\hline \multicolumn{3}{c}{}
\end{longtable}

\subsection*{Macros for model configuration}

Two user-callable macros provided in \texttt{loadmdl.gms}:\begin{description}

\item[\textbf{\texttt{\$macro loaddsys(ff,r)}}] Loads demand system
\texttt{ff} for region \texttt{r} where \texttt{ff} is one of \LES,
\CD or \CDE and \texttt{r} is a region identifier, either a quoted set
element or a subset of the regions \texttt{r}.  This function alters:
\begin{itemize}
\item[\texttt{cd(r)}]  Cobb Douglas demand system flag for model regions
\item[\texttt{les(r)}]  LES demand system flag for model regions
\item[\texttt{cde(r)}] CDE demand system flag
\item[ \texttt{vdfm(),vifm()}] Levels of subsistence and discretionary demand
\item[ \texttt{vom()}] Output level which activate aggregate consumer demand, subsistence and discretionary demand.
\end{itemize}

\item[\texttt{\$macro loadrm(rr)}] Load regions \texttt{rr(r)} as
endogenous elements of the current model.  If \texttt{rr} is a set
which includes all regions in \texttt{r}, then a global multiregional
model is produced.  If \texttt{rr} is singleton or a subset of regions
in \texttt{r}, then a small open economy closure is provided.  This
routine alters: \begin{itemize}

\item[\texttt{rm(r)}]   Set of regions in the model,
\item[\texttt{rx(r)}]   Set of regions in rest of world,
\item[\texttt{rnum(r)}] Set defining the numeraire region -- the region in the model with the largest consumption.
\item[\texttt{vem(i,r)}]        Exports to ROW regions by regions in the model,
\item[\texttt{rtxs\_row(i,r)}]  Average subsidy rates on exports to ROW regions.
\item[\texttt{rowpfx}]          Current account balance for the rest of world regions.
\end{itemize}
\end{description}

\newpage

The user-callable macros provided in \texttt{loadmdl.gms} are invoked
as follows in the \texttt{tgrid.gms} program:

\begin{Verbatim}
*       Load regions in the model, either a global multiregional
*       model containing all regions or a small open economy model
*       with a single region.

$if not set mdl         $set mdl gmr
$if not set rcalc       $set rcalc chn

$if %mdl%==gmr loadrm(r);
$if %mdl%==soe loadrm("%rcalc%");
\end{Verbatim}

In the same program, a loop over functional forms is used to calculate
scenario results for three different demand systems, as shown below.

\subsection*{GDP reporting}

\texttt{gdpcalc.gms} declares the following
identifiers:\footnote{These names may not be used in the calling
program.}

\begin{Verbatim}

set     gdpcat  GDP Categories  /
                expend          Expenditure (C + G + I - (X-M)),
                income          Income (Factor income + taxes),
                valueadded      Sectoral factor earnings plus tax payments,
                total           Total GDP/,

        gdpitem /"X-M",set.g,set.f,
                 revto,revtfd,revtfi,revtf,revtxs,revtms,
                 expend,income,valueadded,chksum/;

alias (gdpitem, gdpi);

parameter       gdp(gdpcat,*,gdpitem)   Real GDP accouting,
                vadd(g, gdpitem, r)     GDP on a value-added basis;
\end{Verbatim}

GDP reporting in \texttt{tgrid.gms} consists of an include statement
at the beginning of the program, just after having read
\texttt{mge.gms}, in a context permitting declarations:

\begin{Verbatim}
*       Read the model and calculate GDP at the benchmark point:

$include %code%mge

$include %code%gdpcalc
\end{Verbatim}

The GDP routine is included following scenarios solutions inside the
loop over alternative demand systems:

\begin{Verbatim}
loop(dsys,

*       Load the demand system:

  loaddsys(dsys,r);

*       Initialize tax instruments at benchmark values:
  rtms(i,s,r) = rtms0(i,s,r);
  rtxs(i,s,r) = rtxs0(i,s,r);

  loop(tlvl,

*       Assign tax rates for the counterfactual simulation:

    rtms(i,s,rcalc)$max(rtms0(i,s,rcalc),0) =
	rtms0(i,s,rcalc) * tlvl.val/averate;
    rtxs(i,rcalc,s)$max(-rtxs0(i,rcalc,s),0) =
	rtxs0(i,rcalc,s) * tlvl.val/averate;

*       Compute the equilibrium values:

$include gtap12.gen
    solve gtap12 using mcp;
    abort$(gtap12.objval>1e-3) "Simulation fails: gtap12.";


*       Store the GDP results for this simulation:

$include %code%gdpcalc

    gdpresults(dsys,tlvl,gdpcat,r,gdpitem) = gdp(gdpcat,r,gdpitem);

...

))

\end{Verbatim}
